\color{unidarkgreen}{\section[(обязательное) Уставки функций РЗА]{(обязательное)\\Уставки функций РЗА}\label{app:set}}
\color{black}

%%%%%%%%%%%%%%%%%%%%%%%%%%%%%%%%%%%%%%%%%%%%%%%%%%%%%%%%%% ДЗ %%%%%%%%%%%%%%%%%%%%%%%%%%%%%%%%%%%%%%%%%%%%%%%%%%%%%%%%%%%%%%%%%%
\par
Параметры для настройки ДЗ приведены в таблице \ref{dz:tbl0}.

\footnotesize
\begin{longtable}{|>{\centering\arraybackslash}m{5.3cm}|>{\centering\arraybackslash}m{3.3cm}|>{\centering\arraybackslash}m{4.2cm}|>{\centering\arraybackslash}m{1.8cm}|>{\centering\arraybackslash}m{1cm}|}
\caption{Общие параметры для настройки ДЗ\hfill\vspace{-0.5\baselineskip}}\label{dz:tbl0}\\ 
\hline
\rowcolor{gray!30}
Параметр (Параметр на ИЧМ) & Условное обозначение на схеме & Значение/ Диапазон & Единица измерения & Шаг \\ 
\hline
\endfirsthead
\caption*{\hspace{3pt}\emph{Продолжение таблицы \ref{dz:tbl0}\hfill\vspace{-0.5\baselineskip}}} \\ % сделано по ГОСТ 2.105 п.6.8.7
\hline
\rowcolor{gray!30}
Параметр (Параметр на ИЧМ) & Условное обозначение на схеме & Значение/ Диапазон & Единица измерения & Шаг \\ 
%\hline
\endhead
%\hline
\endfoot
\hline
\endlastfoot
%==+t1*LLN0|SV_LVDIS
\multicolumn{5}{|c|}{ ДЗ-фф 1 ст } \\ \hline 
\centering Ввод функции в работу (Ввод\_функции)  & \centering SGF1 & \centering Не предусмотрено \\ Предусмотрено & \centering -- & \centering \arraybackslash -- \\
\hline
\centering Угол наклона верхней стороны характеристики срабатывания для компенсации нагрузки (Фкн)  & \centering -- & \centering -45,0...0,0 & \centering град. & \centering \arraybackslash 0,1 \\
\hline
\centering Активное сопротивление срабатывания (Rср)  & \centering -- & \centering 0,01...500,00 & \centering Ом & \centering \arraybackslash 0,01 \\
\hline
\centering Реактивное сопротивление срабатывания (Xср)  & \centering -- & \centering 0,01...500,00 & \centering Ом & \centering \arraybackslash 0,01 \\
\hline
\centering Угол наклона характеристики срабатывания (Фхс)  & \centering -- & \centering 30,0...89,9 & \centering град. & \centering \arraybackslash 0,1 \\
\hline
\centering Режим направленности (Направленность)  & \centering SGF3 & \centering Ненаправленная \\ Прямонаправленная \\ Обратнонаправленная & \centering -- & \centering \arraybackslash -- \\
\hline
\centering Подхват от ненаправленного пуска или отдельной ступени (Подхват)  & \centering SGF9 & \centering Не предусмотрено \\ Предусмотрено & \centering -- & \centering \arraybackslash -- \\
\hline
\centering Режим контроля от БК (Реж\_БК)  & \centering SGF7 & \centering С контролем от БК-б \\ С контролем от БК-м & \centering -- & \centering \arraybackslash -- \\
\hline
\centering Режим контроля от БНН (Реж\_БНН)  & \centering SGF8 & \centering Без контроля \\ С контролем & \centering -- & \centering \arraybackslash -- \\
\hline
\centering Режим контроля от ОПО (Реж\_ОПО)  & \centering SGF6 & \centering По току \\ По току и напряжению \\ По току или по току и напряжению & \centering -- & \centering \arraybackslash -- \\
\hline
\centering Тип ПО (Контр\_ПО)  & \centering SGF5 & \centering Без контроля \\ С контролем от БК \\ С контролем от ОПО & \centering -- & \centering \arraybackslash -- \\
\hline
\centering Вывод направленности при АУ (Ненапр\_АУ)  & \centering SGF4 & \centering Не предусмотрено \\ Предусмотрено & \centering -- & \centering \arraybackslash -- \\
\hline
\centering Выдержка времени срабатывания (Tср)  & \centering T1 & \centering 0...10000 & \centering мс & \centering \arraybackslash 10 \\
\hline
\multicolumn{5}{|c|}{ ДЗ-фф 2 ст } \\ \hline 
\centering Ввод функции в работу (Ввод\_функции)  & \centering SGF1 & \centering Не предусмотрено \\ Предусмотрено & \centering -- & \centering \arraybackslash -- \\
\hline
\centering Угол наклона верхней стороны характеристики срабатывания для компенсации нагрузки (Фкн)  & \centering -- & \centering -45,0...0,0 & \centering град. & \centering \arraybackslash 0,1 \\
\hline
\centering Активное сопротивление срабатывания (Rср)  & \centering -- & \centering 0,01...500,00 & \centering Ом & \centering \arraybackslash 0,01 \\
\hline
\centering Реактивное сопротивление срабатывания (Xср)  & \centering -- & \centering 0,01...500,00 & \centering Ом & \centering \arraybackslash 0,01 \\
\hline
\centering Угол наклона характеристики срабатывания (Фхс)  & \centering -- & \centering 30,0...89,9 & \centering град. & \centering \arraybackslash 0,1 \\
\hline
\centering Режим направленности (Направленность)  & \centering SGF3 & \centering Ненаправленная \\ Прямонаправленная \\ Обратнонаправленная & \centering -- & \centering \arraybackslash -- \\
\hline
\centering Режим контроля от БК (Реж\_БК)  & \centering SGF7 & \centering С контролем от БК-б \\ С контролем от БК-м & \centering -- & \centering \arraybackslash -- \\
\hline
\centering Режим контроля от БНН (Реж\_БНН)  & \centering SGF8 & \centering Без контроля \\ С контролем & \centering -- & \centering \arraybackslash -- \\
\hline
\centering Режим контроля от ОПО (Реж\_ОПО)  & \centering SGF6 & \centering По току \\ По току и напряжению \\ По току или по току и напряжению & \centering -- & \centering \arraybackslash -- \\
\hline
\centering Тип ПО (Контр\_ПО)  & \centering SGF5 & \centering Без контроля \\ С контролем от БК \\ С контролем от ОПО & \centering -- & \centering \arraybackslash -- \\
\hline
\centering Вывод направленности при АУ (Ненапр\_АУ)  & \centering SGF4 & \centering Не предусмотрено \\ Предусмотрено & \centering -- & \centering \arraybackslash -- \\
\hline
\centering Выдержка времени срабатывания (Tср)  & \centering T1 & \centering 0...10000 & \centering мс & \centering \arraybackslash 10 \\
\hline
\multicolumn{5}{|c|}{ ДЗ-фф 3 ст } \\ \hline 
\centering Ввод функции в работу (Ввод\_функции)  & \centering SGF1 & \centering Не предусмотрено \\ Предусмотрено & \centering -- & \centering \arraybackslash -- \\
\hline
\centering Угол наклона верхней стороны характеристики срабатывания для компенсации нагрузки (Фкн)  & \centering -- & \centering -45,0...0,0 & \centering град. & \centering \arraybackslash 0,1 \\
\hline
\centering Активное сопротивление срабатывания (Rср)  & \centering -- & \centering 0,01...500,00 & \centering Ом & \centering \arraybackslash 0,01 \\
\hline
\centering Реактивное сопротивление срабатывания (Xср)  & \centering -- & \centering 0,01...500,00 & \centering Ом & \centering \arraybackslash 0,01 \\
\hline
\centering Угол наклона характеристики срабатывания (Фхс)  & \centering -- & \centering 30,0...89,9 & \centering град. & \centering \arraybackslash 0,1 \\
\hline
\centering Режим направленности (Направленность)  & \centering SGF3 & \centering Ненаправленная \\ Прямонаправленная \\ Обратнонаправленная & \centering -- & \centering \arraybackslash -- \\
\hline
\centering Подхват от ненаправленного пуска или отдельной ступени (Подхват)  & \centering SGF9 & \centering Не предусмотрено \\ Предусмотрено & \centering -- & \centering \arraybackslash -- \\
\hline
\centering Режим контроля от БК (Реж\_БК)  & \centering SGF7 & \centering С контролем от БК-б \\ С контролем от БК-м & \centering -- & \centering \arraybackslash -- \\
\hline
\centering Режим контроля от БНН (Реж\_БНН)  & \centering SGF8 & \centering Без контроля \\ С контролем & \centering -- & \centering \arraybackslash -- \\
\hline
\centering Режим контроля от ОПО (Реж\_ОПО)  & \centering SGF6 & \centering По току \\ По току и напряжению \\ По току или по току и напряжению & \centering -- & \centering \arraybackslash -- \\
\hline
\centering Тип ПО (Контр\_ПО)  & \centering SGF5 & \centering Без контроля \\ С контролем от БК \\ С контролем от ОПО & \centering -- & \centering \arraybackslash -- \\
\hline
\centering Вывод направленности при АУ (Ненапр\_АУ)  & \centering SGF4 & \centering Не предусмотрено \\ Предусмотрено & \centering -- & \centering \arraybackslash -- \\
\hline
\centering Выдержка времени срабатывания (Tср)  & \centering T1 & \centering 0...10000 & \centering мс & \centering \arraybackslash 10 \\
\hline
\multicolumn{5}{|c|}{ ДЗ-фф 4 ст } \\ \hline 
\centering Ввод функции в работу (Ввод\_функции)  & \centering SGF1 & \centering Не предусмотрено \\ Предусмотрено & \centering -- & \centering \arraybackslash -- \\
\hline
\centering Угол наклона верхней стороны характеристики срабатывания для компенсации нагрузки (Фкн)  & \centering -- & \centering -45,0...0,0 & \centering град. & \centering \arraybackslash 0,1 \\
\hline
\centering Активное сопротивление срабатывания (Rср)  & \centering -- & \centering 0,01...500,00 & \centering Ом & \centering \arraybackslash 0,01 \\
\hline
\centering Реактивное сопротивление срабатывания (Xср)  & \centering -- & \centering 0,01...500,00 & \centering Ом & \centering \arraybackslash 0,01 \\
\hline
\centering Угол наклона характеристики срабатывания (Фхс)  & \centering -- & \centering 30,0...89,9 & \centering град. & \centering \arraybackslash 0,1 \\
\hline
\centering Режим направленности (Направленность)  & \centering SGF3 & \centering Ненаправленная \\ Прямонаправленная \\ Обратнонаправленная & \centering -- & \centering \arraybackslash -- \\
\hline
\centering Режим контроля от БК (Реж\_БК)  & \centering SGF7 & \centering С контролем от БК-б \\ С контролем от БК-м & \centering -- & \centering \arraybackslash -- \\
\hline
\centering Режим контроля от БНН (Реж\_БНН)  & \centering SGF8 & \centering Без контроля \\ С контролем & \centering -- & \centering \arraybackslash -- \\
\hline
\centering Режим контроля от ОПО (Реж\_ОПО)  & \centering SGF6 & \centering По току \\ По току и напряжению \\ По току или по току и напряжению & \centering -- & \centering \arraybackslash -- \\
\hline
\centering Тип ПО (Контр\_ПО)  & \centering SGF5 & \centering Без контроля \\ С контролем от БК \\ С контролем от ОПО & \centering -- & \centering \arraybackslash -- \\
\hline
\centering Вывод направленности при АУ (Ненапр\_АУ)  & \centering SGF4 & \centering Не предусмотрено \\ Предусмотрено & \centering -- & \centering \arraybackslash -- \\
\hline
\centering Выдержка времени срабатывания (Tср)  & \centering T1 & \centering 0...10000 & \centering мс & \centering \arraybackslash 10 \\
\hline
\multicolumn{5}{|c|}{ АУ ДЗ } \\ \hline 
\centering Выбор ускоряемой ступени (Выбор\_ст)  & \centering SGF1 & \centering 1 ступень \\ 2 ступень \\ 3 ступень \\ 4 ступень \\ 1 или 3 ступень \\ 2 или 4 ступень & \centering -- & \centering \arraybackslash -- \\
\hline
\multicolumn{5}{|c|}{ ОПО } \\ \hline 
\centering Ток срабатывания ПО ДЗ по току (Iср)  & \centering I\verb|>>| & \centering 0,10...5,00 & \centering о.е. & \centering \arraybackslash 0,01 \\
\hline
\centering Ток срабатывания ПО ДЗ по току и напряжению (IсрU)  & \centering I> & \centering 0,10...5,00 & \centering о.е. & \centering \arraybackslash 0,01 \\
\hline
\centering Напряжение срабатывания ПО по току и напряжению (Uср)  & \centering U< & \centering 2,0...100,0 & \centering В & \centering \arraybackslash 0,1 \\
\hline
\centering Режим контроля от БНН (Реж\_БНН)  & \centering SGF1 & \centering Не предусмотрено \\ Предусмотрено & \centering -- & \centering \arraybackslash -- \\
\hline
\multicolumn{5}{|c|}{ ОН } \\ \hline 
\centering Угол направленности во II квадранте (Фнапр\_II)  & \centering φ2 & \centering 90,0...130,0 & \centering град. & \centering \arraybackslash 0,1 \\
\hline
\centering Угол направленности в IV квадранте (Фнапр\_IV)  & \centering φ1 & \centering -40,0...0,0 & \centering град. & \centering \arraybackslash 0,1 \\
\hline
\multicolumn{5}{|c|}{ ОУ ДЗ } \\ \hline 
\centering Выбор оперативно ускоряемой ступени (Выбор\_ст\_ОУ)  & \centering SGF1 & \centering 1 ступень \\ 2 ступень \\ 3 ступень \\ 4 ступень \\ 1 или 3 ступень \\ 2 или 4 ступень & \centering -- & \centering \arraybackslash -- \\
\hline
\centering Выдержка времени срабатывания оперативного ускорения (T\_ОУ)  & \centering T1 & \centering 0...3000 & \centering мс & \centering \arraybackslash 10 \\
\hline
\multicolumn{5}{|c|}{ ДЗ } \\ \hline 
\centering Учет нагрузочного режима в характеристике срабатывания (Вырез\_нагрузки)  & \centering SGF2 & \centering Не предусмотрено \\ Предусмотрено & \centering -- & \centering \arraybackslash -- \\
\hline
\centering Активное сопротивление для отстройки от режима максимальной нагрузки (Rнг)  & \centering -- & \centering 0,01...500,00 & \centering Ом & \centering \arraybackslash 0,01 \\
\hline
\centering Угол наклона для отстройки от режима максимальной нагрузки (Фнг)  & \centering -- & \centering 5,0...60,0 & \centering град. & \centering \arraybackslash 0,1 \\
\hline
%===t1
\end{longtable}
\normalsize



%%%%%%%%%%%%%%%%%%%%%%%%%%%%%%%%%%%%%%%%%%%%%%%%%%%%%%%%%% БК  %%%%%%%%%%%%%%%%%%%%%%%%%%%%%%%%%%%%%%%%%%%%%%%%%%%%%%%%%%%%%%%%%%
\par
В таблице \ref{bk:tbl1} приведены параметры, необходимые для настройки БК.

\footnotesize
\begin{longtable}{|>{\centering\arraybackslash}m{5.3cm}|>{\centering\arraybackslash}m{3.3cm}|>{\centering\arraybackslash}m{4.2cm}|>{\centering\arraybackslash}m{1.8cm}|>{\centering\arraybackslash}m{1cm}|}
\caption{Параметры для настройки БК\hfill\vspace{-0.5\baselineskip}}\label{bk:tbl1}\\ 
\hline
\rowcolor{gray!20}
Параметр (Параметр на ИЧМ) & Условное обозначение на схеме & Значение/ Диапазон & Единица измерения & Шаг \\ 
\hline
\endfirsthead
\caption*{\hspace{3pt}\emph{Продолжение таблицы \ref{bk:tbl1}\hfill\vspace{-0.5\baselineskip}}} \\ % сделано по ГОСТ 2.105 п.6.8.7
\hline
\rowcolor{gray!20}
Параметр (Параметр на ИЧМ) & Условное обозначение на схеме & Значение/ Диапазон & Единица измерения & Шаг \\ 
%\hline
\endhead
%\hline
\endfoot
%\hline
\endlastfoot
%==+t1*RDIS1|LVPSB
\centering Ввод функции в работу (Ввод\_функции)  & \centering SGF1 & \centering Не предусмотрено \\ Предусмотрено & \centering -- & \centering \arraybackslash -- \\
\hline
\centering Приращение тока ПП чувствительного органа (dI1\_чув)  & \centering dI1/dt> & \centering 0,08...3,00 & \centering о.е. & \centering \arraybackslash 0,01 \\
\hline
\centering Приращение тока ПП грубого органа (dI1\_груб)  & \centering dI1/dt\verb|>>| & \centering 0,12...5,00 & \centering о.е. & \centering \arraybackslash 0,01 \\
\hline
\centering Приращение тока ОП чувствительного органа (dI2\_чув)  & \centering dI2/dt> & \centering 0,04...1,50 & \centering о.е. & \centering \arraybackslash 0,01 \\
\hline
\centering Приращение тока ОП грубого органа (dI2\_груб)  & \centering dI2/dt\verb|>>| & \centering 0,06...2,50 & \centering о.е. & \centering \arraybackslash 0,01 \\
\hline
\centering Время ввода быстродействующих ступеней ДЗ от чувствительных органов (Тб\_чув)  & \centering T1 & \centering 200...1000 & \centering мс & \centering \arraybackslash 10 \\
\hline
\centering Время ввода быстродействующих ступеней ДЗ от грубых органов (Тб\_груб)  & \centering T2 & \centering 200...1000 & \centering мс & \centering \arraybackslash 10 \\
\hline
\centering Время ввода медленнодействующих ступеней ДЗ (Тм)  & \centering T3 & \centering 1000...15000 & \centering мс & \centering \arraybackslash 10 \\
\hline
\centering Ускоренный возврат логики БК при отключении выключателя (Ускор\_возврат)  & \centering SGF2 & \centering Не предусмотрено \\ Предусмотрено & \centering -- & \centering \arraybackslash -- \\
\hline
%===t1
\end{longtable}
\normalsize



%%%%%%%%%%%%%%%%%%%%%%%%%%%%%%%%%%%%%%%%%%%%%%%%%%%%%%%%%% АУ %%%%%%%%%%%%%%%%%%%%%%%%%%%%%%%%%%%%%%%%%%%%%%%%%%%%%%%%%%%%%%%%%%
\par
В таблице \ref{au:tbl1} приведены параметры, необходимые для настройки АУ.

\footnotesize
\begin{longtable}{|>{\centering\arraybackslash}m{5.3cm}|>{\centering\arraybackslash}m{3.3cm}|>{\centering\arraybackslash}m{4.2cm}|>{\centering\arraybackslash}m{1.8cm}|>{\centering\arraybackslash}m{1cm}|}
\caption{Параметры для настройки АУ\hfill\vspace{-0.5\baselineskip}}\label{au:tbl1}\\ 
\hline
\rowcolor{gray!20}
Параметр (Параметр на ИЧМ) & Условное обозначение на схеме & Значение/ Диапазон & Единица измерения & Шаг \\ 
\hline
\endfirsthead
\caption*{\hspace{3pt}\emph{Продолжение таблицы \ref{au:tbl1}\hfill\vspace{-0.5\baselineskip}}} \\ % сделано по ГОСТ 2.105 п.6.8.7
\hline
\rowcolor{gray!20}
Параметр (Параметр на ИЧМ) & Условное обозначение на схеме & Значение/ Диапазон & Единица измерения & Шаг \\ 
%\hline
\endhead
%\hline
\endfoot
%\hline
\endlastfoot
%==+t1*PSOF1|LVLINAUA
\centering Ввод функции в работу (Ввод\_функции)  & \centering SGF1 & \centering Не предусмотрено \\ Предусмотрено & \centering -- & \centering \arraybackslash -- \\
\hline
\centering Автоматическое ускорение МТЗ (АУ\_МТЗ)  & \centering SGF2 & \centering Не предусмотрено \\ Предусмотрено & \centering -- & \centering \arraybackslash -- \\
\hline
\centering Автоматическое ускорение ДЗ (АУ\_ДЗ)  & \centering SGF3 & \centering Не предусмотрено \\ Предусмотрено & \centering -- & \centering \arraybackslash -- \\
\hline
\centering Автоматическое ускорение ТЗНП (АУ\_ТЗНП)  & \centering SGF4 & \centering Не предусмотрено \\ Предусмотрено & \centering -- & \centering \arraybackslash -- \\
\hline
\centering Выдержка времени срабатывания ускоряемой ступени МТЗ (Тср\_МТЗ)  & \centering T1 & \centering 0...3000 & \centering мс & \centering \arraybackslash 10 \\
\hline
\centering Выдержка времени срабатывания ускоряемой ступени ДЗ (Тср\_ДЗ)  & \centering T2 & \centering 0...3000 & \centering мс & \centering \arraybackslash 10 \\
\hline
\centering Выдержка времени срабатывания ускоряемой ступени ТЗНП (Тср\_ТЗНП)  & \centering T3 & \centering 0...3000 & \centering мс & \centering \arraybackslash 10 \\
\hline
\centering Время ввода ускорения при включении выключателя (Твв)  & \centering T4 & \centering 0...5000 & \centering мс & \centering \arraybackslash 10 \\
\hline
%===t1
\end{longtable}
\normalsize


%%%%%%%%%%%%%%%%%%%%%%%%%%%%%%%%%%%%%%%%%%%%%%%%%%%%%%%%%% МТЗ %%%%%%%%%%%%%%%%%%%%%%%%%%%%%%%%%%%%%%%%%%%%%%%%%%%%%%%%%%%%%%%%%% \footnotesize
\par
В таблице \ref{mtz:tbl1} приведены параметры, необходимые для настройки МТЗ.

\footnotesize
\begin{longtable}{|>{\centering\arraybackslash}m{5.3cm}|>{\centering\arraybackslash}m{3.3cm}|>{\centering\arraybackslash}m{4.2cm}|>{\centering\arraybackslash}m{1.8cm}|>{\centering\arraybackslash}m{1cm}|}
\caption{Параметры для настройки ступени МТЗ\hfill\vspace{-0.5\baselineskip}}\label{mtz:tbl1}\\ 
\hline
\rowcolor{gray!30}
Параметр (Параметр на ИЧМ) & Условное обозначение на схеме & Значение/ Диапазон & Единица измерения & Шаг \\ 
\hline
\endfirsthead
\caption*{\hspace{3pt}\emph{Продолжение таблицы \ref{mtz:tbl1}\hfill\vspace{-0.5\baselineskip}}} \\ % сделано по ГОСТ 2.105 п.6.8.7
\hline
\rowcolor{gray!30}
Параметр (Параметр на ИЧМ) & Условное обозначение на схеме & Значение/ Диапазон & Единица измерения & Шаг \\ 
%\hline
\endhead
%\hline
\endfoot
\hline
\endlastfoot
%==+t1*PTOC1|SV_LVTUVTOC
\multicolumn{5}{|c|}{ МТЗ 1 ст } \\ \hline 
\centering Ввод функции в работу (Ввод\_функции)  & \centering SGF1 & \centering Не предусмотрено \\ Предусмотрено & \centering -- & \centering \arraybackslash -- \\
\hline
\centering Ток срабатывания (Iср)  & \centering I> & \centering 0,10...30,00 & \centering о.е. & \centering \arraybackslash 0,01 \\
\hline
\centering Ток срабатывания грубого органа в режиме управляющего напряжения (Iср\_груб.)  & \centering -- & \centering 0,10...30,00 & \centering о.е. & \centering \arraybackslash 0,01 \\
\hline
\centering Ток срабатывания функции включения под нагрузку (Iфвн)  & \centering -- & \centering 0,10...30,00 & \centering о.е. & \centering \arraybackslash 0,01 \\
\hline
\centering Ввод функции включения на "холодную" нагрузку (Ввод\_ФВН)  & \centering SGF2 & \centering Не предусмотрено \\ Предусмотрено & \centering -- & \centering \arraybackslash -- \\
\hline
\centering Время отключенного состояния выключателя (Тзагруб.)  & \centering T2 & \centering 0...20000 & \centering мс & \centering \arraybackslash 10 \\
\hline
\centering Время ввода загрубления (Тсброс.загруб.)  & \centering T3 & \centering 0...20000 & \centering мс & \centering \arraybackslash 10 \\
\hline
\centering Тип пуска по напряжению (Тип\_КПОН)  & \centering SGF6 & \centering Управляющее напряжение \\ Вольтметровая блокировка & \centering -- & \centering \arraybackslash -- \\
\hline
\centering Режим контроля от БНТ (Реж\_БНТ)  & \centering SGF7 & \centering Не предусмотрено \\ Предусмотрено & \centering -- & \centering \arraybackslash -- \\
\hline
\centering Режим направленности (Направленность)  & \centering SGF3 & \centering Ненаправленная \\ Прямонаправленная \\ Обратнонаправленная & \centering -- & \centering \arraybackslash -- \\
\hline
\centering Режим ОНМ при неисправности ЦН (Реж\_БНН\_ОНМ)  & \centering SGF4 & \centering Блокировка \\ Деблокировка & \centering -- & \centering \arraybackslash -- \\
\hline
\centering Вывод направленности при АУ (Ненапр\_АУ)  & \centering SGF5 & \centering Не предусмотрено \\ Предусмотрено & \centering -- & \centering \arraybackslash -- \\
\hline
\centering Режим КПОН при неисправности ЦН (Реж\_БНН\_КПОН)  & \centering SGF8 & \centering Деблокировка \\ Блокировка & \centering -- & \centering \arraybackslash -- \\
\hline
\centering Режим контроля от КПОН (Реж\_КПОН)  & \centering SGF9 & \centering Не предусмотрено \\ Предусмотрено & \centering -- & \centering \arraybackslash -- \\
\hline
\centering Характеристика срабатывания (Характеристика)  & \centering -- & \centering Кривая F \\ Кривая E \\ Кривая D \\ Кривая B \\ Кривая A \\ Кривая C \\ Независимая \\ Кривая RXIDG & \centering -- & \centering \arraybackslash -- \\
\hline
\centering Коэффициент регулировки времени срабатывания зависимой характеристики (К)  & \centering -- & \centering 0,05...15,00 & \centering -- & \centering \arraybackslash 0,01 \\
\hline
\centering Независимая выдержка времени срабатывания (Tср)  & \centering T1 & \centering 0...100000 & \centering мс & \centering \arraybackslash 10 \\
\hline
\multicolumn{5}{|c|}{ МТЗ 2 ст } \\ \hline 
\centering Ввод функции в работу (Ввод\_функции)  & \centering SGF1 & \centering Не предусмотрено \\ Предусмотрено & \centering -- & \centering \arraybackslash -- \\
\hline
\centering Ток срабатывания (Iср)  & \centering I> & \centering 0,10...30,00 & \centering о.е. & \centering \arraybackslash 0,01 \\
\hline
\centering Ток срабатывания грубого органа в режиме управляющего напряжения (Iср\_груб.)  & \centering -- & \centering 0,10...30,00 & \centering о.е. & \centering \arraybackslash 0,01 \\
\hline
\centering Ток срабатывания функции включения под нагрузку (Iфвн)  & \centering -- & \centering 0,10...30,00 & \centering о.е. & \centering \arraybackslash 0,01 \\
\hline
\centering Ввод функции включения на "холодную" нагрузку (Ввод\_ФВН)  & \centering SGF2 & \centering Не предусмотрено \\ Предусмотрено & \centering -- & \centering \arraybackslash -- \\
\hline
\centering Время отключенного состояния выключателя (Тзагруб.)  & \centering T2 & \centering 0...20000 & \centering мс & \centering \arraybackslash 10 \\
\hline
\centering Время ввода загрубления (Тсброс.загруб.)  & \centering T3 & \centering 0...20000 & \centering мс & \centering \arraybackslash 10 \\
\hline
\centering Тип пуска по напряжению (Тип\_КПОН)  & \centering SGF6 & \centering Управляющее напряжение \\ Вольтметровая блокировка & \centering -- & \centering \arraybackslash -- \\
\hline
\centering Режим контроля от БНТ (Реж\_БНТ)  & \centering SGF7 & \centering Не предусмотрено \\ Предусмотрено & \centering -- & \centering \arraybackslash -- \\
\hline
\centering Режим направленности (Направленность)  & \centering SGF3 & \centering Ненаправленная \\ Прямонаправленная \\ Обратнонаправленная & \centering -- & \centering \arraybackslash -- \\
\hline
\centering Режим ОНМ при неисправности ЦН (Реж\_БНН\_ОНМ)  & \centering SGF4 & \centering Блокировка \\ Деблокировка & \centering -- & \centering \arraybackslash -- \\
\hline
\centering Вывод направленности при АУ (Ненапр\_АУ)  & \centering SGF5 & \centering Не предусмотрено \\ Предусмотрено & \centering -- & \centering \arraybackslash -- \\
\hline
\centering Режим КПОН при неисправности ЦН (Реж\_БНН\_КПОН)  & \centering SGF8 & \centering Деблокировка \\ Блокировка & \centering -- & \centering \arraybackslash -- \\
\hline
\centering Режим контроля от КПОН (Реж\_КПОН)  & \centering SGF9 & \centering Не предусмотрено \\ Предусмотрено & \centering -- & \centering \arraybackslash -- \\
\hline
\centering Характеристика срабатывания (Характеристика)  & \centering -- & \centering Кривая F \\ Кривая E \\ Кривая D \\ Кривая B \\ Кривая A \\ Кривая C \\ Независимая \\ Кривая RXIDG & \centering -- & \centering \arraybackslash -- \\
\hline
\centering Коэффициент регулировки времени срабатывания зависимой характеристики (К)  & \centering -- & \centering 0,05...15,00 & \centering -- & \centering \arraybackslash 0,01 \\
\hline
\centering Независимая выдержка времени срабатывания (Tср)  & \centering T1 & \centering 0...100000 & \centering мс & \centering \arraybackslash 10 \\
\hline
\multicolumn{5}{|c|}{ МТЗ 3 ст } \\ \hline 
\centering Ввод функции в работу (Ввод\_функции)  & \centering SGF1 & \centering Не предусмотрено \\ Предусмотрено & \centering -- & \centering \arraybackslash -- \\
\hline
\centering Ток срабатывания (Iср)  & \centering I> & \centering 0,10...30,00 & \centering о.е. & \centering \arraybackslash 0,01 \\
\hline
\centering Ток срабатывания грубого органа в режиме управляющего напряжения (Iср\_груб.)  & \centering -- & \centering 0,10...30,00 & \centering о.е. & \centering \arraybackslash 0,01 \\
\hline
\centering Ток срабатывания функции включения под нагрузку (Iфвн)  & \centering -- & \centering 0,10...30,00 & \centering о.е. & \centering \arraybackslash 0,01 \\
\hline
\centering Ввод функции включения на "холодную" нагрузку (Ввод\_ФВН)  & \centering SGF2 & \centering Не предусмотрено \\ Предусмотрено & \centering -- & \centering \arraybackslash -- \\
\hline
\centering Время отключенного состояния выключателя (Тзагруб.)  & \centering T2 & \centering 0...20000 & \centering мс & \centering \arraybackslash 10 \\
\hline
\centering Время ввода загрубления (Тсброс.загруб.)  & \centering T3 & \centering 0...20000 & \centering мс & \centering \arraybackslash 10 \\
\hline
\centering Тип пуска по напряжению (Тип\_КПОН)  & \centering SGF6 & \centering Управляющее напряжение \\ Вольтметровая блокировка & \centering -- & \centering \arraybackslash -- \\
\hline
\centering Режим контроля от БНТ (Реж\_БНТ)  & \centering SGF7 & \centering Не предусмотрено \\ Предусмотрено & \centering -- & \centering \arraybackslash -- \\
\hline
\centering Режим направленности (Направленность)  & \centering SGF3 & \centering Ненаправленная \\ Прямонаправленная \\ Обратнонаправленная & \centering -- & \centering \arraybackslash -- \\
\hline
\centering Режим ОНМ при неисправности ЦН (Реж\_БНН\_ОНМ)  & \centering SGF4 & \centering Блокировка \\ Деблокировка & \centering -- & \centering \arraybackslash -- \\
\hline
\centering Вывод направленности при АУ (Ненапр\_АУ)  & \centering SGF5 & \centering Не предусмотрено \\ Предусмотрено & \centering -- & \centering \arraybackslash -- \\
\hline
\centering Режим КПОН при неисправности ЦН (Реж\_БНН\_КПОН)  & \centering SGF8 & \centering Деблокировка \\ Блокировка & \centering -- & \centering \arraybackslash -- \\
\hline
\centering Режим контроля от КПОН (Реж\_КПОН)  & \centering SGF9 & \centering Не предусмотрено \\ Предусмотрено & \centering -- & \centering \arraybackslash -- \\
\hline
\centering Характеристика срабатывания (Характеристика)  & \centering -- & \centering Кривая F \\ Кривая E \\ Кривая D \\ Кривая B \\ Кривая A \\ Кривая C \\ Независимая \\ Кривая RXIDG & \centering -- & \centering \arraybackslash -- \\
\hline
\centering Коэффициент регулировки времени срабатывания зависимой характеристики (К)  & \centering -- & \centering 0,05...15,00 & \centering -- & \centering \arraybackslash 0,01 \\
\hline
\centering Независимая выдержка времени срабатывания (Tср)  & \centering T1 & \centering 0...100000 & \centering мс & \centering \arraybackslash 10 \\
\hline
\multicolumn{5}{|c|}{ МТЗ 4 ст } \\ \hline 
\centering Ввод функции в работу (Ввод\_функции)  & \centering SGF1 & \centering Не предусмотрено \\ Предусмотрено & \centering -- & \centering \arraybackslash -- \\
\hline
\centering Ток срабатывания (Iср)  & \centering I> & \centering 0,10...30,00 & \centering о.е. & \centering \arraybackslash 0,01 \\
\hline
\centering Ток срабатывания грубого органа в режиме управляющего напряжения (Iср\_груб.)  & \centering -- & \centering 0,10...30,00 & \centering о.е. & \centering \arraybackslash 0,01 \\
\hline
\centering Ток срабатывания функции включения под нагрузку (Iфвн)  & \centering -- & \centering 0,10...30,00 & \centering о.е. & \centering \arraybackslash 0,01 \\
\hline
\centering Ввод функции включения на "холодную" нагрузку (Ввод\_ФВН)  & \centering SGF2 & \centering Не предусмотрено \\ Предусмотрено & \centering -- & \centering \arraybackslash -- \\
\hline
\centering Время отключенного состояния выключателя (Тзагруб.)  & \centering T2 & \centering 0...20000 & \centering мс & \centering \arraybackslash 10 \\
\hline
\centering Время ввода загрубления (Тсброс.загруб.)  & \centering T3 & \centering 0...20000 & \centering мс & \centering \arraybackslash 10 \\
\hline
\centering Тип пуска по напряжению (Тип\_КПОН)  & \centering SGF6 & \centering Управляющее напряжение \\ Вольтметровая блокировка & \centering -- & \centering \arraybackslash -- \\
\hline
\centering Режим контроля от БНТ (Реж\_БНТ)  & \centering SGF7 & \centering Не предусмотрено \\ Предусмотрено & \centering -- & \centering \arraybackslash -- \\
\hline
\centering Режим направленности (Направленность)  & \centering SGF3 & \centering Ненаправленная \\ Прямонаправленная \\ Обратнонаправленная & \centering -- & \centering \arraybackslash -- \\
\hline
\centering Режим ОНМ при неисправности ЦН (Реж\_БНН\_ОНМ)  & \centering SGF4 & \centering Блокировка \\ Деблокировка & \centering -- & \centering \arraybackslash -- \\
\hline
\centering Вывод направленности при АУ (Ненапр\_АУ)  & \centering SGF5 & \centering Не предусмотрено \\ Предусмотрено & \centering -- & \centering \arraybackslash -- \\
\hline
\centering Режим КПОН при неисправности ЦН (Реж\_БНН\_КПОН)  & \centering SGF8 & \centering Деблокировка \\ Блокировка & \centering -- & \centering \arraybackslash -- \\
\hline
\centering Режим контроля от КПОН (Реж\_КПОН)  & \centering SGF9 & \centering Не предусмотрено \\ Предусмотрено & \centering -- & \centering \arraybackslash -- \\
\hline
\centering Характеристика срабатывания (Характеристика)  & \centering -- & \centering Кривая F \\ Кривая E \\ Кривая D \\ Кривая B \\ Кривая A \\ Кривая C \\ Независимая \\ Кривая RXIDG & \centering -- & \centering \arraybackslash -- \\
\hline
\centering Коэффициент регулировки времени срабатывания зависимой характеристики (К)  & \centering -- & \centering 0,05...15,00 & \centering -- & \centering \arraybackslash 0,01 \\
\hline
\centering Независимая выдержка времени срабатывания (Tср)  & \centering T1 & \centering 0...100000 & \centering мс & \centering \arraybackslash 10 \\
\hline
\multicolumn{5}{|c|}{ КПОН } \\ \hline 
\centering Напряжение срабатывания (Uср)  & \centering U< & \centering 2,0...100,0 & \centering В & \centering \arraybackslash 0,1 \\
\hline
\centering Напряжение срабатывания ОП (U2ср)  & \centering U2> & \centering 2,0...50,0 & \centering В & \centering \arraybackslash 0,1 \\
\hline
\centering Режим пуска (Реж\_пуска)  & \centering SGF1 & \centering По Uмин \\ Комбинированный \\ Внешний & \centering -- & \centering \arraybackslash -- \\
\hline
\multicolumn{5}{|c|}{ ОНМ } \\ \hline 
\centering Угол максимальной чувствительности (Фмч)  & \centering φмч & \centering -179,0...180,0 & \centering град. & \centering \arraybackslash 0,1 \\
\hline
\multicolumn{5}{|c|}{ БНТ } \\ \hline 
\centering Отношение тока второй гармоники к току основной гармоники (Ih2/Ih1)  & \centering I2г/I1г> & \centering 5...100 & \centering \% & \centering \arraybackslash 1 \\
\hline
\centering Ток блокировки (Iср)  & \centering I> & \centering 0,1...15,0 & \centering о.е. & \centering \arraybackslash 0,1 \\
\hline
\centering Перекрестная блокировка (Перекрест\_блок)  & \centering SGF1 & \centering Не предусмотрено \\ Предусмотрено & \centering -- & \centering \arraybackslash -- \\
\hline
\multicolumn{5}{|c|}{ БЛЗШ } \\ \hline 
\centering Выбор ступени блокировки (БлокЛЗШ\_выбор\_ст)  & \centering SGF1 & \centering Не предусмотрено \\ 1 ступень \\ 2 ступень \\ 3 ступень \\ 4 ступень \\ 1 или 3 ступень \\ 2 или 4 ступень & \centering -- & \centering \arraybackslash -- \\
\hline
\multicolumn{5}{|c|}{ АУ МТЗ } \\ \hline 
\centering Выбор ускоряемой ступени (Выбор\_ст)  & \centering SGF1 & \centering 1 ступень \\ 2 ступень \\ 3 ступень \\ 4 ступень \\ 1 или 3 ступень \\ 2 или 4 ступень & \centering -- & \centering \arraybackslash -- \\
\hline
\multicolumn{5}{|c|}{ ОУ МТЗ } \\ \hline 
\centering Выбор оперативно ускоряемой ступени (Выбор\_ст\_ОУ)  & \centering SGF1 & \centering 1 ступень \\ 2 ступень \\ 3 ступень \\ 4 ступень \\ 1 или 3 ступень \\ 2 или 4 ступень & \centering -- & \centering \arraybackslash -- \\
\hline
\centering Выдержка времени срабатывания оперативного ускорения (T\_ОУ)  & \centering T1 & \centering 0...3000 & \centering мс & \centering \arraybackslash 10 \\
\hline
%===t1
\end{longtable}
\normalsize



%%%%%%%%%%%%%%%%%%%%%%%%%%%%%%%%%%%%%%%%%%%%%%%%%%%%%%%%%% ЛЗШ  %%%%%%%%%%%%%%%%%%%%%%%%%%%%%%%%%%%%%%%%%%%%%%%%%%%%%%%%%%%%%%%%%% \footnotesize
\par
В таблице \ref{lzsh:tbl1} приведены параметры, необходимые для настройки ЛЗШ.

\footnotesize
\begin{longtable}{|>{\centering\arraybackslash}m{5.3cm}|>{\centering\arraybackslash}m{3.3cm}|>{\centering\arraybackslash}m{4.2cm}|>{\centering\arraybackslash}m{1.8cm}|>{\centering\arraybackslash}m{1cm}|}
\caption{Параметры для настройки ЛЗШ\hfill\vspace{-0.5\baselineskip}}\label{lzsh:tbl1}\\ 
\hline
\rowcolor{gray!20}
Параметр (Параметр на ИЧМ) & Условное обозначение на схеме & Значение/ Диапазон & Единица измерения & Шаг \\ 
\hline
\endfirsthead
\caption*{\hspace{3pt}\emph{Продолжение таблицы \ref{lzsh:tbl1}\hfill\vspace{-0.5\baselineskip}}} \\ % сделано по ГОСТ 2.105 п.6.8.7
\hline
\rowcolor{gray!20}
Параметр (Параметр на ИЧМ) & Условное обозначение на схеме & Значение/ Диапазон & Единица измерения & Шаг \\ 
%\hline
\endhead
%\hline
\endfoot
%\hline
\endlastfoot
%==+t1*PTOC1|LBPTOC
\centering Ввод функции в работу (Ввод\_функции)  & \centering SGF1 & \centering Не предусмотрено \\ Предусмотрено & \centering -- & \centering \arraybackslash -- \\
\hline
\centering Ток срабатывания (Iср)  & \centering I> & \centering 0,10...30,00 & \centering о.е. & \centering \arraybackslash 0,01 \\
\hline
\centering Ток срабатывания грубого органа в режиме управляющего напряжения (Iср\_груб.)  & \centering -- & \centering 0,10...30,00 & \centering о.е. & \centering \arraybackslash 0,01 \\
\hline
\centering Тип пуска по напряжению (Тип\_КПОН)  & \centering SGF2 & \centering Управляющее напряжение \\ Вольтметровая блокировка & \centering -- & \centering \arraybackslash -- \\
\hline
\centering Режим контроля от БНТ (Реж\_БНТ)  & \centering SGF3 & \centering Не предусмотрено \\ Предусмотрено & \centering -- & \centering \arraybackslash -- \\
\hline
\centering Режим КПОН при неисправности ЦН (Реж\_БНН\_КПОН)  & \centering SGF4 & \centering Деблокировка \\ Блокировка & \centering -- & \centering \arraybackslash -- \\
\hline
\centering Режим контроля от КПОН (Реж\_КПОН)  & \centering SGF5 & \centering Не предусмотрено \\ Предусмотрено & \centering -- & \centering \arraybackslash -- \\
\hline
\centering Выдержка времени неисправности ЛЗШ (Тнеисп)  & \centering T2 & \centering 100...10000 & \centering мс & \centering \arraybackslash 10 \\
\hline
\centering Независимая выдержка времени срабатывания (Tср)  & \centering T1 & \centering 100...10000 & \centering мс & \centering \arraybackslash 10 \\
\hline
%===t1
\end{longtable}
\normalsize


%%%%%%%%%%%%%%%%%%%%%%%%%%%%%%%%%%%%%%%%%%%%%%%%%%%%%%%%%% ЗДЗ  %%%%%%%%%%%%%%%%%%%%%%%%%%%%%%%%%%%%%%%%%%%%%%%%%%%%%%%%%%%%%%%%%% \footnotesize
\par
В таблице \ref{zdz:tbl1} приведены параметры, необходимые для настройки ЗДЗ.

\footnotesize
\begin{longtable}{|>{\centering\arraybackslash}m{5.3cm}|>{\centering\arraybackslash}m{3.3cm}|>{\centering\arraybackslash}m{4.2cm}|>{\centering\arraybackslash}m{1.8cm}|>{\centering\arraybackslash}m{1cm}|}
\caption{Параметры для настройки ЗДЗ\hfill\vspace{-0.5\baselineskip}}\label{zdz:tbl1}\\ 
\hline
\rowcolor{gray!30}
Параметр (Параметр на ИЧМ) & Условное обозначение на схеме & Значение/ Диапазон & Единица измерения & Шаг \\ 
\hline
\endfirsthead
\caption*{\hspace{3pt}\emph{Продолжение таблицы \ref{zdz:tbl1}\hfill\vspace{-0.5\baselineskip}}} \\ % сделано по ГОСТ 2.105 п.6.8.7
\hline
\rowcolor{gray!30}
Параметр (Параметр на ИЧМ) & Условное обозначение на схеме & Значение/ Диапазон & Единица измерения & Шаг \\ 
%\hline
\endhead
%\hline
\endfoot
\hline
\endlastfoot
%==+t1*ARCPTOC1|SV_ARCTOC
\centering Ввод функции в работу (Ввод\_функции)  & \centering SGF1 & \centering Не предусмотрено \\ Предусмотрено & \centering -- & \centering \arraybackslash -- \\
\hline
\centering Ток срабатывания (Iср)  & \centering I> & \centering 0,10...30,00 & \centering о.е. & \centering \arraybackslash 0,01 \\
\hline
\centering Напряжение срабатывания (3U0ср)  & \centering 3U0> & \centering 2,0...150,0 & \centering В & \centering \arraybackslash 0,1 \\
\hline
\centering Блокировка при неисправности ЗДЗ (Блок\_неиспр\_ЗДЗ)  & \centering SGF2 & \centering Не предусмотрено \\ Предусмотрено & \centering -- & \centering \arraybackslash -- \\
\hline
\centering Контроль тока (Контр\_тока)  & \centering SGF3 & \centering Не предусмотрено \\ Контроль I \\ Контроль I или 3U0 \\ Внешний & \centering -- & \centering \arraybackslash -- \\
\hline
\centering Режим контроля от БНН (Реж\_БНН)  & \centering SGF4 & \centering Без контроля \\ С контролем & \centering -- & \centering \arraybackslash -- \\
\hline
\centering Выдержка времени неисправности ЗДЗ (Tнеисп)  & \centering T2 & \centering 0...10000 & \centering мс & \centering \arraybackslash 10 \\
\hline
\centering Выдержка времени срабатывания (Tср)  & \centering T1 & \centering 0...10000 & \centering мс & \centering \arraybackslash 10 \\
\hline
%===t1
\end{longtable}
\normalsize


%%%%%%%%%%%%%%%%%%%%%%%%%%%%%%%%%%%%%%%%%%%%%%%%%%%%%%%%%% ТК ЗДЗ %%%%%%%%%%%%%%%%%%%%%%%%%%%%%%%%%%%%%%%%%%%%%%%%%%%%%%%%%%%%%%%%%%

В таблице \ref{tkzdz:tbl1} приведены параметры, необходимые для настройки ТК ЗДЗ.

\footnotesize
\begin{longtable}{|>{\centering\arraybackslash}m{5.3cm}|>{\centering\arraybackslash}m{3.3cm}|>{\centering\arraybackslash}m{4.2cm}|>{\centering\arraybackslash}m{1.8cm}|>{\centering\arraybackslash}m{1cm}|}
\caption{Параметры для настройки ТК ЗДЗ\hfill\vspace{-0.5\baselineskip}}\label{tkzdz:tbl1}\\ 
\hline
\rowcolor{gray!20}
Параметр (Параметр на ИЧМ) & Условное обозначение на схеме & Значение/ Диапазон & Единица измерения & Шаг \\ 
\hline
\endfirsthead
%\caption{\emph{Продолжение\hfill\vspace{-0.5\baselineskip}}} \\ % Отступ обозначения таблицы от таблицы и выравнивание надписи слева
\caption*{\hspace{3pt}\emph{Продолжение таблицы \ref{tkzdz:tbl1}\hfill\vspace{-0.5\baselineskip}}} \\ % сделано по ГОСТ 2.105 п.6.8.7
\hline
\rowcolor{gray!20}
Параметр (Параметр на ИЧМ) & Условное обозначение на схеме & Значение/ Диапазон & Единица измерения & Шаг \\ 
%\hline
\endhead
%\hline
\endfoot
%\hline
\endlastfoot
%==+t1*PTOC1|LVARCTOC
\centering Ввод функции в работу (Ввод\_функции)  & \centering SGF1 & \centering Не предусмотрено \\ Предусмотрено & \centering -- & \centering \arraybackslash -- \\
\hline
\centering Ток срабатывания (Iср)  & \centering I> & \centering 0,10...30,00 & \centering о.е. & \centering \arraybackslash 0,01 \\
\hline
%===t1
\end{longtable}
\normalsize



%%%%%%%%%%%%%%%%%%%%%%%%%%%%%%%%%%%%%%%%%%%%%%%%%%%%%%%%%% ЗОП %%%%%%%%%%%%%%%%%%%%%%%%%%%%%%%%%%%%%%%%%%%%%%%%%%%%%%%%%%%%%%%%%%
\par
В таблице \ref{zop:tbl1} приведены параметры, необходимые для настройки ЗОП.

\footnotesize
\begin{longtable}{|>{\centering\arraybackslash}m{5.3cm}|>{\centering\arraybackslash}m{3.3cm}|>{\centering\arraybackslash}m{4.2cm}|>{\centering\arraybackslash}m{1.8cm}|>{\centering\arraybackslash}m{1cm}|}
\caption{Параметры для настройки ЗОП\hfill\vspace{-0.5\baselineskip}}\label{zop:tbl1}\\ 
\hline
\rowcolor{gray!20}
Параметр (Параметр на ИЧМ) & Условное обозначение на схеме & Значение/ Диапазон & Единица измерения & Шаг \\ 
\hline
\endfirsthead
\caption*{\hspace{3pt}\emph{Продолжение таблицы \ref{zop:tbl1}\hfill\vspace{-0.5\baselineskip}}} \\ % сделано по ГОСТ 2.105 п.6.8.7
\hline
\rowcolor{gray!20}
Параметр (Параметр на ИЧМ) & Условное обозначение на схеме & Значение/ Диапазон & Единица измерения & Шаг \\ 
%\hline
\endhead
%\hline
\endfoot
%\hline
\endlastfoot
%==+t1*NSPTOC1|LVNSTOC
\centering Ввод функции в работу (Ввод\_функции)  & \centering SGF1 & \centering Не предусмотрено \\ Предусмотрено & \centering -- & \centering \arraybackslash -- \\
\hline
\centering Ток срабатывания ОП (I2ср)  & \centering I2> & \centering 0,04...4,00 & \centering о.е. & \centering \arraybackslash 0,01 \\
\hline
\centering Отношение тока ОП к току ПП (Kнесим)  & \centering I2/I1> & \centering 0,02...1,00 & \centering о.е. & \centering \arraybackslash 0,01 \\
\hline
\centering Режим работы ЗОП (Реж\_ЗОП)  & \centering SGF2 & \centering по I2 \\ по I2/I1 \\ по I2 или по I2/I1 & \centering -- & \centering \arraybackslash -- \\
\hline
\centering Выдержка времени срабатывания (Tср)  & \centering T1 & \centering 100...20000 & \centering мс & \centering \arraybackslash 10 \\
\hline
%===t1
\end{longtable}
\normalsize


%%%%%%%%%%%%%%%%%%%%%%%%%%%%%%%%%%%%%%%%%%%%%%%%%%%%%%%%%% ТЗНП %%%%%%%%%%%%%%%%%%%%%%%%%%%%%%%%%%%%%%%%%%%%%%%%%%%%%%%%%%%%%%%%%% \footnotesize
\par

В таблице \ref{tznp:tbl1} приведены параметры, необходимые для настройки ТЗНП.

\footnotesize
\begin{longtable}{|>{\centering\arraybackslash}m{5.3cm}|>{\centering\arraybackslash}m{3.3cm}|>{\centering\arraybackslash}m{4.2cm}|>{\centering\arraybackslash}m{1.8cm}|>{\centering\arraybackslash}m{1cm}|}
\caption{Параметры для настройки ступеней ТЗНП\hfill\vspace{-0.5\baselineskip}}\label{tznp:tbl1}\\ 
\hline
\rowcolor{gray!30}
Параметр (Параметр на ИЧМ) & Условное обозначение на схеме & Значение/ Диапазон & Единица измерения & Шаг \\ 
\hline
\endfirsthead
\caption*{\hspace{3pt}\emph{Продолжение таблицы \ref{tznp:tbl1}\hfill\vspace{-0.5\baselineskip}}} \\ % сделано по ГОСТ 2.105 п.6.8.7
\hline
\rowcolor{gray!30}
Параметр (Параметр на ИЧМ) & Условное обозначение на схеме & Значение/ Диапазон & Единица измерения & Шаг \\ 
%\hline
\endhead
%\hline
\endfoot
\hline
\endlastfoot
%==+t1*ZSPTOC1|SV_LVZSTOC
\multicolumn{5}{|c|}{ ТЗНП 1 ст } \\ \hline 
\centering Ввод функции в работу (Ввод\_функции)  & \centering SGF1 & \centering Не предусмотрено \\ Предусмотрено & \centering -- & \centering \arraybackslash -- \\
\hline
\centering Режим оперативного блокирования ЧС (Реж\_БЧС)  & \centering SGF8 & \centering Не предусмотрено \\ Предусмотрено & \centering -- & \centering \arraybackslash -- \\
\hline
\centering Режим контроля от БНТ (Реж\_БНТ)  & \centering SGF7 & \centering Не предусмотрено \\ Предусмотрено & \centering -- & \centering \arraybackslash -- \\
\hline
\centering Ток срабатывания НП (3I0ср)  & \centering 3I0> & \centering 0,10...30,00 & \centering о.е. & \centering \arraybackslash 0,01 \\
\hline
\centering Вывод направленности при АУ (Ненапр\_АУ)  & \centering SGF4 & \centering Не предусмотрено \\ Предусмотрено & \centering -- & \centering \arraybackslash -- \\
\hline
\centering Режим направленности (Направленность)  & \centering SGF2 & \centering Ненаправленная \\ Прямонаправленная \\ Обратнонаправленная & \centering -- & \centering \arraybackslash -- \\
\hline
\centering Режим ОНМ НП при неисправности ЦН (Реж\_БНН\_ОНМ\_НП)  & \centering SGF3 & \centering Блокировка \\ Деблокировка & \centering -- & \centering \arraybackslash -- \\
\hline
\centering Режим ПО 3U0 при неисправности ЦН (Реж\_БНН\_3U0)  & \centering SGF5 & \centering Блокировка \\ Деблокировка & \centering -- & \centering \arraybackslash -- \\
\hline
\centering Напряжение срабатывания НП (3U0ср)  & \centering 3U0> & \centering 2,0...100,0 & \centering В & \centering \arraybackslash 0,1 \\
\hline
\centering Режим контроля от ПО 3U0 (Реж\_3U0)  & \centering SGF6 & \centering Не предусмотрено \\ Предусмотрено & \centering -- & \centering \arraybackslash -- \\
\hline
\centering Характеристика срабатывания (Характеристика)  & \centering -- & \centering Кривая F \\ Кривая E \\ Кривая D \\ Кривая B \\ Кривая A \\ Кривая C \\ Независимая \\ Кривая RXIDG & \centering -- & \centering \arraybackslash -- \\
\hline
\centering Коэффициент регулировки времени срабатывания зависимой характеристики (К)  & \centering -- & \centering 0,05...15,00 & \centering -- & \centering \arraybackslash 0,01 \\
\hline
\centering Независимая выдержка времени срабатывания (Tср)  & \centering T1 & \centering 0...100000 & \centering мс & \centering \arraybackslash 10 \\
\hline
\multicolumn{5}{|c|}{ ТЗНП 2 ст } \\ \hline 
\centering Ввод функции в работу (Ввод\_функции)  & \centering SGF1 & \centering Не предусмотрено \\ Предусмотрено & \centering -- & \centering \arraybackslash -- \\
\hline
\centering Режим оперативного блокирования ЧС (Реж\_БЧС)  & \centering SGF8 & \centering Не предусмотрено \\ Предусмотрено & \centering -- & \centering \arraybackslash -- \\
\hline
\centering Режим контроля от БНТ (Реж\_БНТ)  & \centering SGF7 & \centering Не предусмотрено \\ Предусмотрено & \centering -- & \centering \arraybackslash -- \\
\hline
\centering Ток срабатывания НП (3I0ср)  & \centering 3I0> & \centering 0,10...30,00 & \centering о.е. & \centering \arraybackslash 0,01 \\
\hline
\centering Вывод направленности при АУ (Ненапр\_АУ)  & \centering SGF4 & \centering Не предусмотрено \\ Предусмотрено & \centering -- & \centering \arraybackslash -- \\
\hline
\centering Режим направленности (Направленность)  & \centering SGF2 & \centering Ненаправленная \\ Прямонаправленная \\ Обратнонаправленная & \centering -- & \centering \arraybackslash -- \\
\hline
\centering Режим ОНМ НП при неисправности ЦН (Реж\_БНН\_ОНМ\_НП)  & \centering SGF3 & \centering Блокировка \\ Деблокировка & \centering -- & \centering \arraybackslash -- \\
\hline
\centering Режим ПО 3U0 при неисправности ЦН (Реж\_БНН\_3U0)  & \centering SGF5 & \centering Блокировка \\ Деблокировка & \centering -- & \centering \arraybackslash -- \\
\hline
\centering Напряжение срабатывания НП (3U0ср)  & \centering 3U0> & \centering 2,0...100,0 & \centering В & \centering \arraybackslash 0,1 \\
\hline
\centering Режим контроля от ПО 3U0 (Реж\_3U0)  & \centering SGF6 & \centering Не предусмотрено \\ Предусмотрено & \centering -- & \centering \arraybackslash -- \\
\hline
\centering Характеристика срабатывания (Характеристика)  & \centering -- & \centering Кривая F \\ Кривая E \\ Кривая D \\ Кривая B \\ Кривая A \\ Кривая C \\ Независимая \\ Кривая RXIDG & \centering -- & \centering \arraybackslash -- \\
\hline
\centering Коэффициент регулировки времени срабатывания зависимой характеристики (К)  & \centering -- & \centering 0,05...15,00 & \centering -- & \centering \arraybackslash 0,01 \\
\hline
\centering Независимая выдержка времени срабатывания (Tср)  & \centering T1 & \centering 0...100000 & \centering мс & \centering \arraybackslash 10 \\
\hline
\multicolumn{5}{|c|}{ ТЗНП 3 ст } \\ \hline 
\centering Ввод функции в работу (Ввод\_функции)  & \centering SGF1 & \centering Не предусмотрено \\ Предусмотрено & \centering -- & \centering \arraybackslash -- \\
\hline
\centering Режим оперативного блокирования ЧС (Реж\_БЧС)  & \centering SGF8 & \centering Не предусмотрено \\ Предусмотрено & \centering -- & \centering \arraybackslash -- \\
\hline
\centering Режим контроля от БНТ (Реж\_БНТ)  & \centering SGF7 & \centering Не предусмотрено \\ Предусмотрено & \centering -- & \centering \arraybackslash -- \\
\hline
\centering Ток срабатывания НП (3I0ср)  & \centering 3I0> & \centering 0,10...30,00 & \centering о.е. & \centering \arraybackslash 0,01 \\
\hline
\centering Вывод направленности при АУ (Ненапр\_АУ)  & \centering SGF4 & \centering Не предусмотрено \\ Предусмотрено & \centering -- & \centering \arraybackslash -- \\
\hline
\centering Режим направленности (Направленность)  & \centering SGF2 & \centering Ненаправленная \\ Прямонаправленная \\ Обратнонаправленная & \centering -- & \centering \arraybackslash -- \\
\hline
\centering Режим ОНМ НП при неисправности ЦН (Реж\_БНН\_ОНМ\_НП)  & \centering SGF3 & \centering Блокировка \\ Деблокировка & \centering -- & \centering \arraybackslash -- \\
\hline
\centering Режим ПО 3U0 при неисправности ЦН (Реж\_БНН\_3U0)  & \centering SGF5 & \centering Блокировка \\ Деблокировка & \centering -- & \centering \arraybackslash -- \\
\hline
\centering Напряжение срабатывания НП (3U0ср)  & \centering 3U0> & \centering 2,0...100,0 & \centering В & \centering \arraybackslash 0,1 \\
\hline
\centering Режим контроля от ПО 3U0 (Реж\_3U0)  & \centering SGF6 & \centering Не предусмотрено \\ Предусмотрено & \centering -- & \centering \arraybackslash -- \\
\hline
\centering Характеристика срабатывания (Характеристика)  & \centering -- & \centering Кривая F \\ Кривая E \\ Кривая D \\ Кривая B \\ Кривая A \\ Кривая C \\ Независимая \\ Кривая RXIDG & \centering -- & \centering \arraybackslash -- \\
\hline
\centering Коэффициент регулировки времени срабатывания зависимой характеристики (К)  & \centering -- & \centering 0,05...15,00 & \centering -- & \centering \arraybackslash 0,01 \\
\hline
\centering Независимая выдержка времени срабатывания (Tср)  & \centering T1 & \centering 0...100000 & \centering мс & \centering \arraybackslash 10 \\
\hline
\multicolumn{5}{|c|}{ ТЗНП 4 ст } \\ \hline 
\centering Ввод функции в работу (Ввод\_функции)  & \centering SGF1 & \centering Не предусмотрено \\ Предусмотрено & \centering -- & \centering \arraybackslash -- \\
\hline
\centering Режим оперативного блокирования ЧС (Реж\_БЧС)  & \centering SGF8 & \centering Не предусмотрено \\ Предусмотрено & \centering -- & \centering \arraybackslash -- \\
\hline
\centering Режим контроля от БНТ (Реж\_БНТ)  & \centering SGF7 & \centering Не предусмотрено \\ Предусмотрено & \centering -- & \centering \arraybackslash -- \\
\hline
\centering Ток срабатывания НП (3I0ср)  & \centering 3I0> & \centering 0,10...30,00 & \centering о.е. & \centering \arraybackslash 0,01 \\
\hline
\centering Вывод направленности при АУ (Ненапр\_АУ)  & \centering SGF4 & \centering Не предусмотрено \\ Предусмотрено & \centering -- & \centering \arraybackslash -- \\
\hline
\centering Режим направленности (Направленность)  & \centering SGF2 & \centering Ненаправленная \\ Прямонаправленная \\ Обратнонаправленная & \centering -- & \centering \arraybackslash -- \\
\hline
\centering Режим ОНМ НП при неисправности ЦН (Реж\_БНН\_ОНМ\_НП)  & \centering SGF3 & \centering Блокировка \\ Деблокировка & \centering -- & \centering \arraybackslash -- \\
\hline
\centering Режим ПО 3U0 при неисправности ЦН (Реж\_БНН\_3U0)  & \centering SGF5 & \centering Блокировка \\ Деблокировка & \centering -- & \centering \arraybackslash -- \\
\hline
\centering Напряжение срабатывания НП (3U0ср)  & \centering 3U0> & \centering 2,0...100,0 & \centering В & \centering \arraybackslash 0,1 \\
\hline
\centering Режим контроля от ПО 3U0 (Реж\_3U0)  & \centering SGF6 & \centering Не предусмотрено \\ Предусмотрено & \centering -- & \centering \arraybackslash -- \\
\hline
\centering Характеристика срабатывания (Характеристика)  & \centering -- & \centering Кривая F \\ Кривая E \\ Кривая D \\ Кривая B \\ Кривая A \\ Кривая C \\ Независимая \\ Кривая RXIDG & \centering -- & \centering \arraybackslash -- \\
\hline
\centering Коэффициент регулировки времени срабатывания зависимой характеристики (К)  & \centering -- & \centering 0,05...15,00 & \centering -- & \centering \arraybackslash 0,01 \\
\hline
\centering Независимая выдержка времени срабатывания (Tср)  & \centering T1 & \centering 0...100000 & \centering мс & \centering \arraybackslash 10 \\
\hline
\multicolumn{5}{|c|}{ АУ ТЗНП } \\ \hline 
\centering Выбор ускоряемой ступени (Выбор\_ст)  & \centering SGF1 & \centering 1 ступень \\ 2 ступень \\ 3 ступень \\ 4 ступень \\ 1 или 3 ступень \\ 2 или 4 ступень & \centering -- & \centering \arraybackslash -- \\
\hline
\multicolumn{5}{|c|}{ ОНМ НП } \\ \hline 
\centering Угол максимальной чувствительности (Фмч)  & \centering φмч & \centering -179,0...180,0 & \centering град. & \centering \arraybackslash 0,1 \\
\hline
\multicolumn{5}{|c|}{ БНТ НП } \\ \hline 
\centering Отношение тока второй гармоники к току основной гармоники (Ih2/Ih1)  & \centering I2г/I1г> & \centering 5...100 & \centering \% & \centering \arraybackslash 1 \\
\hline
\centering Ток блокировки НП (3I0ср)  & \centering 3I0> & \centering 0,1...15,0 & \centering о.е. & \centering \arraybackslash 0,1 \\
\hline
\multicolumn{5}{|c|}{ ОУ ТЗНП } \\ \hline 
\centering Выбор оперативно ускоряемой ступени (Выбор\_ст\_ОУ)  & \centering SGF1 & \centering 1 ступень \\ 2 ступень \\ 3 ступень \\ 4 ступень \\ 1 или 3 ступень \\ 2 или 4 ступень & \centering -- & \centering \arraybackslash -- \\
\hline
\centering Выдержка времени срабатывания оперативного ускорения (T\_ОУ)  & \centering T1 & \centering 0...3000 & \centering мс & \centering \arraybackslash 10 \\
\hline
%===t1
\end{longtable}
\normalsize



%%%%%%%%%%%%%%%%%%%%%%%%%%%%%%%%%%%%%%%%%%%%%%%%%%%%%%%%%% УРОВ %%%%%%%%%%%%%%%%%%%%%%%%%%%%%%%%%%%%%%%%%%%%%%%%%%%%%%%%%%%%%%%%%%
\par
В таблице \ref{urov35:tbl1} приведены параметры, необходимые для настройки УРОВ.

\footnotesize
\begin{longtable}{|>{\centering\arraybackslash}m{5.3cm}|>{\centering\arraybackslash}m{3.3cm}|>{\centering\arraybackslash}m{4.2cm}|>{\centering\arraybackslash}m{1.8cm}|>{\centering\arraybackslash}m{1cm}|}
\caption{Параметры для настройки УРОВ\hfill\vspace{-0.5\baselineskip}}\label{urov35:tbl1}\\ 
\hline
\rowcolor{gray!20}
Параметр (Параметр на ИЧМ) & Условное обозначение на схеме & Значение/ Диапазон & Единица измерения & Шаг \\ 
\hline
\endfirsthead
\caption*{\hspace{3pt}\emph{Продолжение таблицы \ref{urov35:tbl1}\hfill\vspace{-0.5\baselineskip}}} \\ % сделано по ГОСТ 2.105 п.6.8.7
\hline
\rowcolor{gray!20}
Параметр (Параметр на ИЧМ) & Условное обозначение на схеме & Значение/ Диапазон & Единица измерения & Шаг \\ 
%\hline
\endhead
%\hline
\endfoot
%\hline
\endlastfoot
%==+t1*GENRBRF1|SV_TPBRF
\centering Ввод функции в работу (Ввод\_функции)  & \centering SGF1 & \centering Не предусмотрено \\ Предусмотрено & \centering -- & \centering \arraybackslash -- \\
\hline
\centering Ток срабатывания (Iср)  & \centering I> & \centering 0,05...0,50 & \centering о.е. & \centering \arraybackslash 0,01 \\
\hline
\centering Ускорение при блокировке отключения В (Уск\_при\_блк\_откл\_В)  & \centering SGF2 & \centering Не предусмотрено \\ Предусмотрено & \centering -- & \centering \arraybackslash -- \\
\hline
\centering УРОВ с подхватом по току (Подхват\_по\_току)  & \centering SGF3 & \centering Не предусмотрено \\ Предусмотрено & \centering -- & \centering \arraybackslash -- \\
\hline
\centering Контроль по току при действии "на себя" (Контр\_тока\_на\_себя)  & \centering SGF6 & \centering Не предусмотрено \\ Предусмотрено по внутр. ПО \\ Предусмотрено по внеш. ПО & \centering -- & \centering \arraybackslash -- \\
\hline
\centering Контроль ЭМО (Контроль\_ЭМО)  & \centering SGF4 & \centering Не предусмотрено \\ Предусмотрено по ЭМО1 \\ Предусмотрено по ЭМО1 и ЭМО2 & \centering -- & \centering \arraybackslash -- \\
\hline
\centering Действие внешнего УРОВ на вышестоящий выключатель (Действ\_на\_выш\_выкл)  & \centering SGF5 & \centering Не предусмотрено \\ Предусмотрено & \centering -- & \centering \arraybackslash -- \\
\hline
\centering Выдержка времени срабатывания (Tср)  & \centering T1 & \centering 0...1000 & \centering мс & \centering \arraybackslash 10 \\
\hline
%===t1
\end{longtable}
\normalsize